\section{Implementation results}
\label{sec:results}

\subsection{Training behavior}
\label{sec:training-behavior}
\reporttodo{Describe the loss and validation recall curves for each trained
model: the selected epoch, signs of overfitting or underfitting, and training
time. Record implementation issues found during training and how they were
fixed.}

\reportfigure{Images/a3_learning_curves.png}{Training loss and validation
recall for each trained model on Flickr8k.}{fig:learning-curves}

\subsection{Quantitative results}
\label{sec:quantitative}
\reporttodo{Report Recall@1, Recall@5, Recall@10 and median rank in both
directions for every model on the test split, using the checkpoint selected on
validation. Give parameter counts, training time and inference cost. State the
number of seeds, or give mean and standard deviation if more seeds are run.
Make every reported number traceable (Handbook Section 4.2): name the config
file, split files, checkpoint, log or experiment ID, and the commit or git tag
that produced it.}

\begin{table}[htbp]
\centering\small
\caption{Flickr8k test retrieval over 1,000 images and 5,000 captions, with the
checkpoint selected on validation. R@$K$ in percent; MedR is the median rank
(lower is better).}
\label{tab:retrieval-results}
\setlength{\tabcolsep}{3.5pt}
\begin{tabular}{lrrrrrrrr}
\toprule
& \multicolumn{4}{c}{Image to text} & \multicolumn{4}{c}{Text to image}\\
\cmidrule(lr){2-5}\cmidrule(lr){6-9}
Model & R@1 & R@5 & R@10 & MedR & R@1 & R@5 & R@10 & MedR\\
\midrule
Image-only baseline & \pending & \pending & \pending & \pending & \pending & \pending & \pending & \pending\\
Text-only baseline & \pending & \pending & \pending & \pending & \pending & \pending & \pending & \pending\\
Simple fusion & \pending & \pending & \pending & \pending & \pending & \pending & \pending & \pending\\
Improved model & \pending & \pending & \pending & \pending & \pending & \pending & \pending & \pending\\
\bottomrule
\end{tabular}
\end{table}


\subsection{Compute cost}
\label{sec:compute}
\reporttodo{Compare parameters, training time, time to encode the test set,
retrieval latency, and peak memory. Name the hardware and software versions for
each run. If runs used different hardware, say why their times cannot be
compared directly.}

\begin{table}[htbp]
\centering\small
\caption{Compute cost. Training time is measured on the hardware listed;
test-set encoding is measured on \pending{} at batch size \pending.}
\label{tab:compute-cost}
\begin{tabular}{lrrrrl}
\toprule
& & Train time & Test encoding & Peak memory & Training\\
Model & Parameters & (s) & (s) & (GB) & hardware\\
\midrule
Image-only baseline & \pending & \pending & \pending & \pending & \pending\\
Text-only baseline & \pending & \pending & \pending & \pending & \pending\\
Simple fusion & \pending & \pending & \pending & \pending & \pending\\
Improved model & \pending & \pending & \pending & \pending & \pending\\
\bottomrule
\end{tabular}
\end{table}


\subsection{Qualitative results}
\label{sec:qualitative}
\reporttodo{Show test queries with their top-ranked results from each model:
correct retrievals, difficult but correct cases, and failures. Mark the ground
truth in each ranking. State the rule used to select the examples so the
selection is reproducible.}

\reportfigure{Images/a3_retrieval_examples.png}{Test queries with the top
retrieved results from each model. State the selection rule.}{fig:retrieval-examples}

\subsection{Modality interaction: support and conflict}
\label{sec:interaction}
\reporttodo{Show cases where the two modalities support each other (the fused
model succeeds where a unimodal baseline fails) and cases where they conflict
(the fused model fails where a unimodal baseline succeeds). Count how often
each case occurs on the test split and explain likely causes with evidence
(Handbook Sections 24.1 and 30).}

\reportfigure{Images/a3_interaction_examples.png}{Test cases where the
modalities support each other and where they conflict.}{fig:interaction}

\subsection{Ablation study}
\label{sec:ablation}
\reporttodo{Report the ablation that removes or changes a modality or a fusion
component. Compare the retrieval metrics and state whether the hypothesis is
supported under the decision criterion in \secref{sec:experiments}.}

\begin{table}[htbp]
\centering\small
\caption{Ablation on the Flickr8k \pending{} split. Only the listed component
changes; R@$K$ in percent.}
\label{tab:ablation}
\begin{tabular}{lrrrr}
\toprule
Variant & I2T R@1 & I2T R@10 & T2I R@1 & T2I R@10\\
\midrule
Full model & \pending & \pending & \pending & \pending\\
\pending & \pending & \pending & \pending & \pending\\
\bottomrule
\end{tabular}
\end{table}


\subsection{Error analysis}
\label{sec:errors}
\reporttodo{Group the errors, for example a near-correct caption ranked above
the true one, captions shared by several images, confusion between similar dog
or people scenes, and counting or colour mistakes. For each group give a
description, its frequency or rate, examples, a hypothesized cause, and an
improvement idea. Use the 21 own pairs without a CrowdFlower majority yes and
the expert judgements to separate model errors from metric false negatives.}

\subsection{Comparison and discussion}
\label{sec:discussion}
\reporttodo{Explain which model is better and by how much, in which direction
and on which kinds of query the differences appear, and likely causes in data,
encoders, fusion, loss or training. If the multimodal model does not beat a
unimodal baseline, explain why with evidence (Handbook Section 29.1). Discuss
accuracy, speed and parameter trade-offs. Say which conclusions cannot yet be
drawn. Listing the highest score is not enough.}

\subsection{Limitations and conclusion}
\label{sec:conclusion}
\reporttodo{List data, method and experimental limits (dataset size and topic
bias, metric false negatives, number of seeds, compute). Discuss generalization
and future work. End with a short answer to the research question.}
